\DocumentMetadata{
  pdfversion=2.0,
  pdfstandard=ua-2,
  lang=en-US,
  tagging=on,
  tagging-setup={math/setup=mathml-SE,tabsorder=structure}
}
\documentclass[11pt,letterpaper]{article}
\usepackage[margin=0.68in,includefoot]{geometry}
\usepackage{newcomputermodern}
\usepackage{amsmath,amscd,mathtools}
\usepackage{tikz,adjustbox}
\providecommand{\square}{\mdlgwhtsquare}
\usepackage[hidelinks]{hyperref}
\hypersetup{
  pdftitle={Combinatorics PhD Exam, August 1995},
  pdfauthor={University of Florida Department of Mathematics},
  pdfsubject={Graduate mathematics examination},
  pdfdisplaydoctitle=true,
  bookmarksopen=true
}
\setcounter{secnumdepth}{0}
\setlength{\parindent}{0pt}
\setlength{\parskip}{0.28em}
\setlength{\emergencystretch}{3em}
\setlength{\leftmargini}{2.15em}
\setlength{\leftmarginii}{2.65em}
\setlength{\leftmarginiii}{3em}
\pagestyle{empty}
\raggedbottom
\allowdisplaybreaks
\definecolor{ProblemConcern}{HTML}{7A1F1F}
\newenvironment{problemconcern}{%
  \par\tagstructbegin{tag=Aside}\begingroup\color{ProblemConcern}
}{%
  \par\endgroup\tagstructend\nopagebreak[4]\vspace{0.12em}
}
\NewTaggingSocketPlug{block/list/label}{examproblem}{%
  \tagstructbegin{tag=Lbl}\tagstructend%
  \tagstructbegin{tag=\LBody}%
  \tagstructbegin{tag=\ProblemHeadingTag,title-o={\ProblemHeadingText}}%
  \tagmcbegin{tag=\ProblemHeadingTag,actualtext={\ProblemHeadingText}}%
  #2%
  \tagmcend%
  \tagstructend%
}
\newenvironment{examproblems}[2]{%
  \def\ProblemHeadingTag{#2}%
  \AssignTaggingSocketPlug{block/list/label}{examproblem}%
  \begin{enumerate}%
  \setcounter{enumi}{#1}%
  \renewcommand{\labelenumi}{\textbf{\theenumi.}}%
  \setlength{\topsep}{0.35em}%
  \setlength{\partopsep}{0pt}%
  \setlength{\itemsep}{0.48em}%
  \setlength{\parsep}{0.18em}%
}{\end{enumerate}}
\newenvironment{examsubparts}{%
  \AssignTaggingSocketPlug{block/list/label}{default}%
  \begin{enumerate}%
  \setlength{\topsep}{0.18em}%
  \setlength{\partopsep}{0pt}%
  \setlength{\itemsep}{0.16em}%
  \setlength{\parsep}{0.1em}%
}{\end{enumerate}}
\newenvironment{examinstructions}{%
  \par\begingroup\itshape
}{%
  \par\endgroup\vspace{0.18em}
}
\makeatletter
\renewcommand\subsection{\@startsection{subsection}{2}{0pt}%
  {-1.25ex plus -.25ex minus -.1ex}%
  {0.55ex plus .1ex}%
  {\normalfont\large\bfseries}}
\renewcommand{\@maketitle}{%
  \newpage\null\vskip -1em%
  {\footnotesize\bfseries
    \href{https://www.ufl.edu/}{University of Florida}, \href{https://math.ufl.edu/}{Department of Mathematics}\par}%
  \vskip -0.5em%
  \tagstructbegin{tag=H1,title={Combinatorics PhD Exam, August 1995}}%
  \tagpdfparaOff%
  \tagmcbegin{tag=H1}%
    {\LARGE\bfseries\href{https://gma.math.ufl.edu/exams/combinatorics-phd/}{Combinatorics PhD Exam}, August 1995\par}%
  \tagmcend%
  \tagpdfparaOn%
  \tagstructend%
  \vskip 0.25em%
  \tagmcbegin{artifact}%
    \hrule height 1pt%
  \tagmcend%
  \vskip 1em%
}
\makeatother
\title{Combinatorics PhD Exam, August 1995}
\author{}
\date{}
\begin{document}
\maketitle
\thispagestyle{empty}
\begin{examproblems}{0}{H2}
\def\ProblemHeadingText{Problem 1.}
\item
Every simple planar graph has a vertex of degree at most~\(5\).
\def\ProblemHeadingText{Problem 2.}
\item
This problem has two parts.
\begin{examsubparts}
\item[\textnormal{(a)}]
State the max-flow min-cut theorem for networks.
\item[\textnormal{(b)}]
Show that the max-flow min-cut theorem implies the following version of Menger’s Theorem.

\par
The maximum number of edge disjoint directed paths between two given points \(s,t\) of a directed graph equals the minimum number of edges whose removal separates~\(s\) from~\(t\).
\end{examsubparts}
\def\ProblemHeadingText{Problem 3.}
\item
We have an unlimited supply of bricks of length~\(1\), each painted red, white or blue, and bricks of length~\(2\), each painted green or yellow. Explicitly find \(a_n\) = the number of ways to choose a sequence of bricks of total length~\(n\).
\def\ProblemHeadingText{Problem 4.}
\item
State and prove Dilworth’s Theorem for a finite poset.
\def\ProblemHeadingText{Problem 5.}
\item
Each of~\(n\) gentlemen checks both his hat and his umbrella at a restaurant. Both the hats and the umbrellas are returned randomly. What is the probability that no man gets back both his hat and his umbrella?
\def\ProblemHeadingText{Problem 6.}
\item
\begin{problemconcern}
\textbf{Warning: Suspected error.} The summation upper limit is clearly printed as~\(m\), but~\(m\) is otherwise undefined and no conditions are given under which the displayed identity holds. The source reading has been preserved because the intended correction is not uniquely determined.
\end{problemconcern}
Prove that
\[
\sum_{i=0}^{m} \binom{s}{i}\binom{t}{k-i}=\binom{s+t}{k}.
\]
\def\ProblemHeadingText{Problem 7.}
\item
\(6\) identical black balls and \(n-6\) identical white balls are arranged in a row. How many ways are there to do this so that no~\(3\) consecutive balls are black?
\def\ProblemHeadingText{Problem 8.}
\item
Define a~\(q\)-Hamming code~\(C\) as one whose parity-check matrix~\(H\) has as its columns one non-zero vector from each~\(1\)-dimensional subspace of the vector space \(\mathrm{GF}(q)^r\).
\begin{examsubparts}
\item[\textnormal{(a)}]
What is the length of~\(C\)?
\item[\textnormal{(b)}]
What is the minimum distance in~\(C\)?
\item[\textnormal{(c)}]
Prove that~\(C\) is perfect.
\end{examsubparts}
\def\ProblemHeadingText{Problem 9.}
\item
This problem has two parts.
\begin{examsubparts}
\item[\textnormal{(a)}]
State the Bruck–Ryser–Chowla theorem.
\item[\textnormal{(b)}]
Use this theorem to prove that a projective plane of order~\(6\) cannot exist.
\end{examsubparts}
\end{examproblems}
\end{document}
